\documentclass[11pt]{article}

\usepackage[margin=1in]{geometry}
\usepackage[T1]{fontenc}
\usepackage[utf8]{inputenc}
\usepackage{microtype}
\usepackage{amsmath,amssymb,mathtools,bm}
\usepackage{booktabs,longtable,tabularx,array}
\usepackage{enumitem}
\usepackage{xcolor}
\usepackage{hyperref}
\usepackage[nameinlink,noabbrev]{cleveref}
\usepackage[numbers,sort&compress]{natbib}

\hypersetup{
  colorlinks=true,
  linkcolor=blue!55!black,
  citecolor=green!35!black,
  urlcolor=blue!65!black,
  pdftitle={A Graph Neural Surrogate for Dynamic System-Optimal Traffic Assignment},
  pdfauthor={Research design document}
}

\newcommand{\G}{\mathcal{G}}
\newcommand{\V}{\mathcal{V}}
\newcommand{\E}{\mathcal{E}}
\newcommand{\K}{\mathcal{K}}
\newcommand{\T}{\mathcal{T}}
\newcommand{\R}{\mathbb{R}}
\newcommand{\Eop}{\mathbb{E}}
\newcommand{\ind}{\mathbb{I}}
\newcommand{\relu}{\operatorname{ReLU}}
\newcommand{\softplus}{\operatorname{softplus}}
\newcommand{\softmax}{\operatorname{softmax}}
\newcommand{\argmin}{\operatorname*{arg\,min}}
\newcommand{\poa}{\operatorname{PoA}}
\newcommand{\gap}{\operatorname{Gap}}
\newcommand{\vect}[1]{\bm{#1}}
\newcommand{\trans}{\mathsf{T}}

\title{\textbf{A Graph Neural Surrogate for Dynamic System-Optimal Traffic Assignment}\\
\large Model Design, Training Pipeline, and Reproducible Evaluation Protocol}
\author{Research design document}
\date{August 2026}

\begin{document}
\maketitle

\begin{abstract}
This document specifies how to build and evaluate a graph-neural surrogate for dynamic
system-optimal traffic assignment (DSO-DTA).  The proposed system is not a generic
traffic forecaster.  It is a decision model that maps a road network, current traffic
state, time-varying origin--destination demand, and disruptions to feasible routing
decisions intended to minimize total system travel time.  The recommended design uses
a cell-transmission-model oracle to generate system-optimal demonstrations, a
heterogeneous spatio-temporal graph neural network to represent roads and long-range
origin--destination interactions, and a differentiable feasibility projection to
enforce conservation and capacity constraints.  Supervised imitation is followed by
decision-focused fine-tuning through a differentiable or simulation-based traffic
rollout.  The protocol distinguishes interpolation from genuine generalization by
holding out complete network topologies, demand regimes, and incident patterns.  It
also compares the surrogate with user equilibrium, marginal-cost assignment,
model-predictive control, and the exact oracle.  The resulting study is designed to
test whether expensive repeated DSO optimization can be replaced by fast inference
without hiding infeasibility, instability, unfairness, or simulator mismatch.
\end{abstract}

\noindent\textbf{Status of this document.}
This is a research and implementation blueprint, not a report of completed
experiments.  Numerical thresholds described below are proposed evaluation targets;
they must not be presented as empirical findings until the experiments are run.

\tableofcontents

\section{Research objective and scope}

The motivating phenomenon is the difference between individually rational routing and
the social optimum.  Wardrop's first principle describes a user equilibrium (UE), in
which no traveler can improve their own experienced cost by changing routes, while
Wardrop's second principle describes a system optimum (SO), in which average or total
travel cost is minimized \citep{wardrop1952}.  The two need not coincide.  This
inefficiency is commonly summarized by the price of anarchy, and it can be substantial
in congested transportation models \citep{roughgarden2002,youn2008}.

The proposed research asks:

\begin{quote}
Can a graph neural network approximate the first control action of a dynamic
system-optimal traffic assignment solver on previously unseen demand profiles,
incidents, and road-network topologies, while maintaining physical feasibility and a
small total-travel-time optimality gap?
\end{quote}

The central engineering use case is a navigation platform or traffic-management center
that periodically recommends routes to a controlled population.  The initial paper
should study \emph{routing only}: traffic signals, departure-time control, pricing,
and strategic non-compliance should be left fixed or treated as extensions.  Combining
all of them in the first model would make attribution and theoretical analysis
unnecessarily difficult.

\subsection{Precise intended contribution}

The intended contribution has four parts:

\begin{enumerate}[label=\textbf{C\arabic*:},leftmargin=2.8em]
  \item A receding-horizon DSO-DTA oracle based on the cell transmission model (CTM),
  used to generate training decisions and an upper-performance reference.
  \item A heterogeneous spatio-temporal GNN that represents physical road cells,
  intersection movements, and virtual origin--destination (OD) relations.
  \item A feasibility mechanism---preferably a differentiable projection---that
  enforces nonnegativity, movement capacities, conservation, and admissible turns.
  \item A topology-held-out evaluation that reports decision quality, physical
  feasibility, runtime, robustness, and the fraction of the price-of-anarchy gap
  recovered by the learned controller.
\end{enumerate}

Traffic-assignment surrogates based on graph neural networks already exist.
Prior work has learned static equilibrium flows from OD demand \citep{rahman2023},
incorporated heterogeneous OD links and conservation penalties for static UE
assignment \citep{liu2024hetgat}, and modeled spatio-temporal flows from OD demand
\citep{liu2024hstgsn}.  Consequently, ``using a GNN for traffic assignment'' is not a
sufficient novelty claim.  The defensible gap is the combination of dynamic
\emph{system-optimal control}, hard or explicitly measured feasibility, decision-focused
training, and genuine unseen-topology testing.

\section{Traffic-assignment foundations}

\subsection{Network, commodities, and demand}

Let the directed road network be
\(
  \G=(\V,\E)
\), where intersections or cell boundaries are nodes and directed road segments are
edges.  Let \(\K\) index OD commodities.  Commodity \(k\in\K\) has origin \(o_k\),
destination \(d_k\), and exogenous departure demand \(r_k(t)\) during discrete time
interval \(t\in\T=\{0,\ldots,H-1\}\).  The interval length is \(\Delta t\).

For static exposition, let \(x_e^k\geq 0\) be the flow of commodity \(k\) on link
\(e\), and \(x_e=\sum_k x_e^k\).  The node--arc incidence matrix is
\(B\in\{-1,0,1\}^{|\V|\times|\E|}\).  Static feasibility can be written

\begin{equation}
  B\vect{x}^k = \vect{b}^k,
  \qquad
  \vect{x}^k\geq \vect{0},
  \label{eq:static-conservation}
\end{equation}

where \(b_v^k\) is positive at the origin, negative at the destination, and zero at
intermediate nodes, under the selected incidence convention.

\subsection{Static user equilibrium and system optimum}

A common separable link travel-time model is the Bureau of Public Roads (BPR) form

\begin{equation}
  \tau_e(x_e)
  = \tau_e^0
  \left[
    1+\alpha_e\left(\frac{x_e}{c_e}\right)^{\beta_e}
  \right],
  \label{eq:bpr}
\end{equation}

where \(\tau_e^0\) is free-flow travel time and \(c_e\) is nominal capacity.  Under
standard monotonicity and separability conditions, the Beckmann transformation writes
the UE as the convex program \citep{beckmann1956}

\begin{equation}
  \vect{x}^{\mathrm{UE}}
  \in
  \argmin_{\vect{x}\in\mathcal{F}}
  \sum_{e\in\E}
  \int_0^{x_e}\tau_e(z)\,\mathrm{d}z,
  \label{eq:ue}
\end{equation}

where \(\mathcal{F}\) is the multicommodity feasible-flow set.  The static SO instead
solves

\begin{equation}
  \vect{x}^{\mathrm{SO}}
  \in
  \argmin_{\vect{x}\in\mathcal{F}}
  C(\vect{x}),
  \qquad
  C(\vect{x}) = \sum_{e\in\E} x_e\tau_e(x_e).
  \label{eq:so}
\end{equation}

For differentiable link costs, the marginal social cost is

\begin{equation}
  m_e(x_e)
  = \frac{\mathrm{d}}{\mathrm{d}x_e}
    \big[x_e\tau_e(x_e)\big]
  = \tau_e(x_e)+x_e\tau'_e(x_e).
  \label{eq:marginal-cost}
\end{equation}

This identity motivates one possible neural output: a learned, time-dependent
generalization of the external congestion term \(x_e\tau'_e(x_e)\).  Static UE and SO
can be solved with algorithms based on Frank--Wolfe and related traffic-assignment
methods \citep{frank1956,patriksson1994}.  A neural surrogate is therefore most
valuable in the dynamic, repeatedly solved setting rather than as a replacement for a
single moderate-sized static assignment.

The conventional static price of anarchy is

\begin{equation}
  \poa = \frac{C(\vect{x}^{\mathrm{UE}})}
                   {C(\vect{x}^{\mathrm{SO}})} \geq 1.
  \label{eq:poa}
\end{equation}

For a learned policy \(\pi_\theta\), a more informative normalized quantity is the
fraction of the avoidable gap that remains:

\begin{equation}
  \gap(\pi_\theta)
  = \frac{C(\pi_\theta)-C(\mathrm{SO})}
         {C(\mathrm{UE})-C(\mathrm{SO})+\varepsilon}.
  \label{eq:gap}
\end{equation}

Values near zero indicate that the controller closes most of the UE--SO gap.  Values
above one indicate performance worse than UE.  A value below zero is not credible
unless the policy and oracle were evaluated under different feasible sets, stochastic
realizations, or solver tolerances; such a result should trigger an audit rather than
an optimality claim.

\subsection{Why the dynamic problem is different}

Dynamic traffic assignment adds departure time, queue propagation, storage capacity,
spillback, and time-dependent route interaction.  A route decision at time \(t\) can
change downstream congestion many intervals later.  The dynamic SO problem therefore
cannot generally be reduced to independent shortest-path calculations using current
travel times.  Reviews of DTA distinguish analytical, simulation-based, and
variational-inequality formulations \citep{peeta2001}.

The CTM discretizes the kinematic-wave traffic model into conservation equations that
represent the formation, propagation, and dissipation of queues
\citep{daganzo1994,daganzo1995}.  CTM formulations have also been used to express
single-destination SO-DTA as a linear program \citep{ziliaskopoulos2000}.  This makes
CTM an appropriate initial oracle: it is substantially more realistic than a static
BPR assignment while remaining easier to optimize and differentiate through than a
microscopic simulator.

\section{Dynamic system-optimal oracle}

\subsection{Cell graph and state variables}

Split each sufficiently long road link into cells.  Let \(\mathcal{C}\) be the set of
cells and \(\mathcal{M}\subseteq\mathcal{C}\times\mathcal{C}\) the allowed directed
cell movements.  For cell \(i\), define:

\begin{itemize}[leftmargin=1.5em]
  \item \(N_i\): maximum vehicle storage;
  \item \(F_i\): maximum flow per interval;
  \item \(v_i\): free-flow sending coefficient;
  \item \(w_i\): backward-wave receiving coefficient;
  \item \(n_i^k(t)\): number of commodity-\(k\) vehicles in cell \(i\);
  \item \(y_{ij}^k(t)\): commodity-\(k\) flow from cell \(i\) to \(j\).
\end{itemize}

Total occupancy and movement flow are
\(
  n_i(t)=\sum_k n_i^k(t)
\)
and
\(
  y_{ij}(t)=\sum_k y_{ij}^k(t)
\).
The multicommodity conservation equation is

\begin{equation}
  n_i^k(t+1)
  = n_i^k(t)
  + a_i^k(t)
  + \sum_{j:(j,i)\in\mathcal{M}}y_{ji}^k(t)
  - \sum_{j:(i,j)\in\mathcal{M}}y_{ij}^k(t)
  - z_i^k(t),
  \label{eq:ctm-conservation}
\end{equation}

where \(a_i^k(t)\) is admitted exogenous demand and \(z_i^k(t)\) is destination exit
flow.  If entry demand can queue outside the modeled network, introduce source queues
\(q_k^{\mathrm{src}}(t)\) rather than silently dropping vehicles.

For a triangular fundamental diagram, sending and receiving envelopes can be written

\begin{align}
  S_i(t) &= \min\{v_i n_i(t),F_i(t)\},
  \\
  R_i(t) &= \min\{w_i[N_i-n_i(t)],F_i(t)\}.
  \label{eq:ctm-sr}
\end{align}

An incident changes \(F_i(t)\), \(v_i(t)\), or the set of allowed movements.  At each
interval, total movement flows obey

\begin{align}
  \sum_{j:(i,j)\in\mathcal{M}} y_{ij}(t)
    &\leq S_i(t),
  \\
  \sum_{i:(i,j)\in\mathcal{M}} y_{ij}(t)
    &\leq R_j(t),
  \\
  y_{ij}^k(t)&\geq 0.
  \label{eq:ctm-capacity}
\end{align}

Merge priorities, lane groups, signal phases, and node conflicts can be added as linear
or mixed-integer constraints.  For the first routing-only paper, signal timing should
be held fixed and represented as a known time-dependent movement capacity.

\subsection{Objective}

By discrete-time vehicle accounting, total system travel time over the horizon can be
approximated by accumulated occupancy:

\begin{equation}
  J_{\mathrm{TSTT}}
  = \Delta t\sum_{t=0}^{H-1}
      \left[
        \sum_{i\in\mathcal{C}}n_i(t)
        + \sum_{k\in\K}q_k^{\mathrm{src}}(t)
      \right].
  \label{eq:tstt}
\end{equation}

The source queues must be counted; otherwise a controller can appear effective by
refusing to admit vehicles.  A practical finite-horizon objective is

\begin{equation}
  J = J_{\mathrm{TSTT}}
      + \lambda_{\mathrm{term}}
        \sum_i \omega_i n_i(H)
      + \lambda_{\mathrm{switch}}J_{\mathrm{switch}}
      + \lambda_{\mathrm{fair}}J_{\mathrm{fair}},
  \label{eq:full-objective}
\end{equation}

where the terminal penalty approximates cost beyond the horizon, the switching term
discourages excessive route changes, and the optional fairness term controls disparate
detours.  For the core surrogate paper, set \(\lambda_{\mathrm{fair}}=0\) during the
primary experiments and report fairness as an outcome; adding a fairness objective is
a natural follow-up study.

\subsection{Receding-horizon oracle}

At decision time \(t\), the oracle receives current state \(\vect{s}_t\), a demand
forecast \(\widehat{\vect{r}}_{t:t+H-1}\), and predicted capacities.  It solves

\begin{equation}
  \vect{u}_{t:t+H-1}^{\star}
  \in
  \argmin_{\vect{u},\vect{n},\vect{y}} J
  \quad\text{subject to \cref{eq:ctm-conservation,eq:ctm-capacity} and node constraints.}
  \label{eq:oracle}
\end{equation}

Only the first decision \(\vect{u}_t^\star\) is applied; the problem is re-solved at
the next control interval using new observations.  The GNN should initially learn this
first-action mapping:

\begin{equation}
  \pi^\star:
  (\G,\vect{s}_t,\widehat{\vect{r}}_{t:t+H-1},\vect{c}_{t:t+H-1})
  \mapsto \vect{u}_t^\star.
  \label{eq:first-action-map}
\end{equation}

This is preferable to predicting an entire open-loop horizon because later oracle
actions depend on states that will differ once the surrogate makes an error.

\section{What the neural model should output}

There are three credible output parameterizations.  The first is recommended for the
initial implementation.

\subsection{Recommended: local movement split logits}

For each commodity \(k\), cell or node \(i\), and admissible outgoing movement
\((i,j)\), the network emits a logit \(\ell_{ij}^k(t)\).  A masked softmax gives

\begin{equation}
  p_{ij}^k(t)
  = \frac{\exp(\ell_{ij}^k(t))\ind[(i,j)\in\mathcal{M}_k]}
  {\sum_{h:(i,h)\in\mathcal{M}_k}\exp(\ell_{ih}^k(t))},
  \label{eq:turn-softmax}
\end{equation}

where \(\mathcal{M}_k\) excludes illegal movements and edges from which destination
\(d_k\) is unreachable.  The split proportions determine desired movement flows; a
node allocator or projection reduces them when sending or receiving capacity binds.

Advantages are a variable-size output, local feasibility of the probability simplex,
and compatibility with closed-loop CTM simulation.  The main danger is cycling.
Mitigate it with destination potentials, reachability masks, expected remaining-time
features, and a loop penalty during rollouts.

\subsection{Alternative: learned marginal edge costs}

The model can output a nonnegative externality

\begin{equation}
  \widehat{\xi}_e^k(t)
  = \softplus(g_\theta(\vect{h}_e,\vect{h}_k,\vect{h}_t)),
  \qquad
  \widehat{m}_e^k(t)=\widehat{\tau}_e(t)+\widehat{\xi}_e^k(t).
  \label{eq:learned-cost}
\end{equation}

Time-dependent shortest paths or a differentiable soft shortest-path layer then turn
these costs into route probabilities.  This form is interpretable as a learned dynamic
marginal cost, but dynamic externalities depend on future queue propagation, so they
need not equal the static expression in \cref{eq:marginal-cost}.

\subsection{Alternative: full multicommodity flows}

The network may directly predict \(\widetilde y_{ij}^k(t)\) for an entire horizon.
This makes supervision simple but creates a large tensor and substantial feasibility
burden.  It is useful as an offline speed benchmark, not the preferred deployed
controller.

\section{Heterogeneous spatio-temporal graph architecture}

\subsection{Graph schema}

Use a heterogeneous graph with the entity and relation types in
\cref{tab:heterograph}.  Virtual OD entities are important because an origin and its
destination may be many message-passing hops apart.  Prior traffic-assignment GNNs
have used virtual OD links for this purpose \citep{liu2024hetgat,liu2024hstgsn}.

\begin{table}[ht]
\centering
\caption{Recommended heterogeneous graph schema.}
\label{tab:heterograph}
\begin{tabularx}{\textwidth}{@{}p{0.18\textwidth}p{0.30\textwidth}X@{}}
\toprule
Entity/relation & Examples & Purpose \\
\midrule
Physical cell & occupancy, storage, flow limit, length, speed, lanes & Represents traffic state and local road physics. \\
Intersection & phase, movement conflicts, aggregate pressure & Represents junction-specific constraints. \\
OD token & origin, destination, departure forecast, remaining demand & Represents a variable number of commodities. \\
Cell adjacency & upstream/downstream, turn type, movement capacity & Propagates local congestion and spillback information. \\
OD-to-origin & injection relation & Communicates demand to its entry region. \\
Destination-to-OD & attraction/reachability relation & Communicates destination context over long ranges. \\
Global token & time of day, weather, network-level demand, incident count & Provides scenario context without forcing many graph hops. \\
\bottomrule
\end{tabularx}
\end{table}

\subsection{Features and normalization}

All continuous features should have a physical interpretation and preferably be
normalized by a local scale.  This improves transfer across networks.

\begin{longtable}{@{}p{0.24\textwidth}p{0.69\textwidth}@{}}
\caption{Candidate input features.}\label{tab:features}\\
\toprule
Feature group & Recommended features \\
\midrule
\endfirsthead
\toprule
Feature group & Recommended features \\
\midrule
\endhead
Cell state & \(n_i/N_i\), recent inflow and outflow divided by \(F_i\), speed/free-flow speed, travel-time/free-flow-time, source-queue indicator. \\
Cell geometry & normalized length, lane count, road class, free-flow time, storage, nominal capacity, geographic bearing encoded by sine and cosine. \\
Incident state & capacity multiplier, speed multiplier, closure flag, elapsed incident time, uncertain duration summary. \\
Intersection state & current phase or fixed phase schedule, turn type, conflicting-flow indicator, movement saturation rate. \\
OD state & departures in each forecast interval, unserved demand, origin and destination roles, free-flow OD time, demand relative to nearby capacity. \\
Temporal context & sine/cosine time-of-day, day type, forecast lead time, previous control action. \\
Global context & total active demand, network mean occupancy, fraction of disabled capacity, forecast confidence. \\
\bottomrule
\end{longtable}

Avoid per-network z-score normalization computed from a test topology: it can leak test
statistics and makes deployment dependent on future observations.  Ratios such as
occupancy/storage and flow/capacity are more portable.  Training-set global transforms
may be used for heavy-tailed features such as OD demand.

\subsection{Encoder, message passing, and temporal state}

Each entity type \(r\) receives a separate encoder

\begin{equation}
  \vect{h}_v^{(0)}
  = \phi_r(\vect{x}_v),
  \qquad v\text{ has type }r.
\end{equation}

At spatial layer \(\ell\), relation-specific messages are aggregated:

\begin{align}
  \vect{m}_{u\rightarrow v}^{(\ell,r)}
    &= \psi_r^{(\ell)}
       (\vect{h}_u^{(\ell)},\vect{h}_v^{(\ell)},\vect{e}_{uv}),
  \\
  \vect{m}_v^{(\ell)}
    &= \sum_r\sum_{u\in\mathcal{N}_r(v)}
       \alpha_{uv}^{(\ell,r)}\vect{m}_{u\rightarrow v}^{(\ell,r)},
  \\
  \vect{h}_v^{(\ell+1)}
    &= \operatorname{LayerNorm}
       \left[
         \vect{h}_v^{(\ell)}
         + \rho^{(\ell)}(\vect{m}_v^{(\ell)})
       \right].
  \label{eq:message-passing}
\end{align}

The attention coefficient \(\alpha_{uv}\) may depend on relation, congestion, and OD
context.  Residual connections and normalization reduce optimization difficulty.
Graph attention is not itself the research contribution; it is an implementation
choice motivated by variable neighborhood importance \citep{velickovic2018}.

For temporal information, apply the spatial encoder to each of the last \(L\) observed
states and pass entity embeddings through a gated recurrent unit, temporal convolution,
or causal Transformer.  Begin with a GRU because it is easier to train and cheaper than
a large Transformer.  Forecast OD tokens can use a second temporal encoder over the
future horizon.  Fuse observed state and forecasts using cross-attention or gated
concatenation.

\subsection{Commodity scaling}

A direct \(|\K|\times|\mathcal{M}|\) representation may be too large when every OD pair
is a commodity.  Evaluate three scaling strategies:

\begin{enumerate}[leftmargin=1.5em]
  \item \textbf{Destination batching:} run one destination-conditioned policy for a
  batch of destinations and share the physical graph encoder.
  \item \textbf{OD token sparsification:} instantiate tokens only for OD pairs with
  positive demand in the current horizon.
  \item \textbf{Regional aggregation:} group nearby destinations for upper-level
  guidance, then use conventional local routing inside each region.
\end{enumerate}

The paper should report memory and runtime as both network size and active commodity
count increase.  Hiding commodity growth behind a small benchmark would undermine the
scalability claim.

\section{Hard feasibility and differentiable projection}

A conservation penalty alone does not guarantee a usable routing action.  The preferred
design is to transform neural proposals into the nearest feasible action.  Let
\(\widetilde{\vect{y}}_t\) be proposed movement flows.  Define

\begin{equation}
  \vect{y}_t^{\mathrm{proj}}
  = \argmin_{\vect{y}}
      \frac{1}{2}
      \left\|W_t(\vect{y}-\widetilde{\vect{y}}_t)\right\|_2^2
  \quad\text{subject to}\quad
  A_t\vect{y}=\vect{b}_t,
  \quad
  G_t\vect{y}\leq\vect{h}_t,
  \quad
  \vect{y}\geq\vect{0}.
  \label{eq:projection}
\end{equation}

The equalities encode commodity conservation or desired outflow allocation; the
inequalities encode sending, receiving, and movement limits.  Slack variables with a
large, reported penalty are preferable to infeasibility when source demand exceeds
capacity.  Differentiable optimization layers permit gradients through structured
convex programs \citep{amos2017,agrawal2019}.

Two cautions are essential:

\begin{enumerate}[leftmargin=1.5em]
  \item Projection can dominate runtime.  Report both raw GNN inference and end-to-end
  inference including the projection and route decomposition.
  \item Projection can hide a poor network.  Report the normalized correction
  \(
  \|\vect{y}^{\mathrm{proj}}-\widetilde{\vect{y}}\|/
  (\|\vect{y}^{\mathrm{proj}}\|+\varepsilon)
  \)
  and the raw violation rates.
\end{enumerate}

If the projection is too costly, use a hierarchy: masked softmax for local split
simplexes, an exact CTM node allocator for capacities, and a conservation correction
only at source and sink boundaries.  This is faster but supplies weaker guarantees.

\section{Dataset construction}

\subsection{Networks}

Use a curriculum that combines synthetic and established research networks.  The
Transportation Networks for Research collection includes Braess, Sioux Falls, Anaheim,
Eastern Massachusetts, Chicago, Winnipeg, and other traffic-assignment instances
\citep{transportationNetworks}.  A recommended partition is:

\begin{itemize}[leftmargin=1.5em]
  \item \textbf{Development:} Braess networks and small grids for debugging.
  \item \textbf{Training:} procedurally generated planar networks plus a subset of
  medium benchmark networks.
  \item \textbf{Validation:} entire synthetic graph families and at least one complete
  real benchmark topology excluded from training.
  \item \textbf{Test:} complete topologies never used for model selection, including a
  larger network than any training graph if computationally possible.
\end{itemize}

Split by topology \emph{before} generating demand scenarios.  Randomly splitting
scenarios from the same topology across train and test measures interpolation, not
topological generalization.

\subsection{Dynamic OD demand}

Static TNTP matrices can seed time-varying demand.  For base OD volume \(D_k\), sample

\begin{equation}
  r_k(t)
  = D_k\,s(t;\vect{\eta})\,\zeta_k(t),
  \label{eq:demand}
\end{equation}

where \(s(t;\vect{\eta})\) is a normalized profile with one or more peaks and
\(\zeta_k(t)\) is a correlated multiplicative disturbance.  Include:

\begin{itemize}[leftmargin=1.5em]
  \item low, near-capacity, and oversaturated demand;
  \item narrow and broad peaks;
  \item spatially correlated event demand;
  \item demand shifts toward OD pairs rare or absent in training;
  \item forecast error distinct from realized demand.
\end{itemize}

Record the random seed, complete demand tensor, and forecast supplied to each method.
All compared policies must face the same realized demand.

\subsection{Incident generation}

An incident scenario is described by affected cells, start time, duration, capacity
multiplier, and observation delay.  Separate incident sets into:

\begin{enumerate}[leftmargin=1.5em]
  \item familiar severities at unfamiliar locations;
  \item severities outside the training range;
  \item multiple interacting incidents;
  \item complete closures that change reachability;
  \item sensor-reported incidents with noisy duration forecasts.
\end{enumerate}

Recent robustness studies show that learned and pressure-based traffic controllers can
degrade under incident-driven distribution shift, motivating explicit disruption tests
rather than assuming nominal generalization \citep{nguyen2025trex}.

\subsection{Oracle demonstration generation}

For each scenario:

\begin{enumerate}[leftmargin=1.5em]
  \item Build or load the cell graph and verify reachability for every active OD pair.
  \item Sample realized demand, the forecast visible to the controller, and incidents.
  \item Initialize traffic with a documented warm-up policy.
  \item At each control time, solve the DSO receding-horizon problem.
  \item Store the pre-decision state, forecast, capacities, first oracle action,
  objective value, dual variables when available, and solver diagnostics.
  \item Apply the oracle action to the CTM, advance one control interval, and repeat.
  \item Reject or flag demonstrations with unacceptable primal residual, duality gap,
  or time limit.
\end{enumerate}

Do not treat a timed-out incumbent as exactly optimal.  Store a lower bound when the
solver supplies one, and label each sample with its certified optimality gap.  Samples
may be weighted by solution confidence.

\section{Training objectives}

\subsection{Stage 1: supervised imitation}

Let \(\vect{u}_t^\star\) be the oracle action and
\(\widehat{\vect{u}}_t=\pi_\theta(\cdot)\) the projected model action.  A normalized
Huber or absolute loss is less sensitive to a small number of high-flow movements than
unscaled squared error:

\begin{equation}
  \mathcal{L}_{\mathrm{imit}}
  = \frac{1}{|\mathcal{B}|}
    \sum_{b\in\mathcal{B}}
    \sum_{m,k}
    \omega_{mk}^{(b)}
    \operatorname{Huber}
    \left(
      \frac{\widehat u_{mk}^{(b)}-u_{mk}^{\star(b)}}
           {F_m^{(b)}+\varepsilon}
    \right).
  \label{eq:imitation}
\end{equation}

Flow labels may be non-unique: two different assignments can have almost identical
system cost.  Consequently, imitation error must not be the primary final metric.

\subsection{Physics residuals}

When the model predicts full flows or states, add

\begin{align}
  \mathcal{L}_{\mathrm{cons}}
    &= \sum_{i,k,t}
       \left|
       \widehat n_i^k(t+1)-\widehat n_i^k(t)-a_i^k(t)
       -\sum_j\widehat y_{ji}^k(t)
       +\sum_j\widehat y_{ij}^k(t)
       +\widehat z_i^k(t)
       \right|,
  \\
  \mathcal{L}_{\mathrm{cap}}
    &= \sum_{i,t}
       \relu\!\left(\sum_j\widehat y_{ij}(t)-S_i(t)\right)
       + \sum_{j,t}
       \relu\!\left(\sum_i\widehat y_{ij}(t)-R_j(t)\right),
  \\
  \mathcal{L}_{\mathrm{loop}}
    &= \sum_{k,t}\text{predicted recirculating flow for commodity }k.
  \label{eq:physics-losses}
\end{align}

If hard projection is used, these losses are still useful on raw outputs so that the
network learns not to rely on projection.

\subsection{Stage 2: decision-focused fine-tuning}

Unroll the projected policy through a differentiable CTM or a simulator interface for
\(H_{\mathrm{train}}\) intervals.  Minimize normalized decision regret

\begin{equation}
  \mathcal{L}_{\mathrm{decision}}
  = \frac{J(\pi_\theta)-J(\pi^\star)}
         {|J(\pi^\star)|+\varepsilon}.
  \label{eq:decision-loss}
\end{equation}

The complete objective is

\begin{equation}
  \mathcal{L}
  = \lambda_{\mathrm{imit}}\mathcal{L}_{\mathrm{imit}}
  + \lambda_{\mathrm{dec}}\mathcal{L}_{\mathrm{decision}}
  + \lambda_{\mathrm{cons}}\mathcal{L}_{\mathrm{cons}}
  + \lambda_{\mathrm{cap}}\mathcal{L}_{\mathrm{cap}}
  + \lambda_{\mathrm{loop}}\mathcal{L}_{\mathrm{loop}}
  + \lambda_{\mathrm{proj}}\mathcal{L}_{\mathrm{projection}}.
  \label{eq:total-loss}
\end{equation}

Start with imitation and short rollouts, then gradually increase rollout length and
decision-loss weight.  This curriculum reduces exploding gradients and distribution
shift from compounding policy errors.  If the traffic simulator is nondifferentiable,
use imitation plus reinforcement learning, a learned differentiable dynamics model, or
zeroth-order policy optimization.  A differentiable CTM is preferable for the first
study because it preserves interpretability and reproducibility.

\subsection{On-policy state collection}

Pure behavioral cloning trains on states visited by the oracle, while deployment visits
states caused by the surrogate's own mistakes.  Use iterative data aggregation:

\begin{enumerate}[leftmargin=1.5em]
  \item train on oracle trajectories;
  \item roll out the current surrogate;
  \item query the oracle at selected surrogate-visited states;
  \item add those labeled states to the dataset;
  \item repeat until validation decision regret stabilizes.
\end{enumerate}

Prioritize queries with high uncertainty, large projection corrections, or near-storage
occupancies.  This uses expensive oracle calls where they provide the most information.

\section{From aggregate flows to driver recommendations}

The model controls aggregate flow, but a navigation service displays paths.  The
conversion must be defined explicitly.

\subsection{Route-set method}

For each active OD pair and departure interval:

\begin{enumerate}[leftmargin=1.5em]
  \item Generate \(K_{\mathrm{path}}\) reasonable loop-free candidate paths using
  free-flow time, current time, and learned social costs.
  \item Form the path--movement incidence matrix \(A_k\).
  \item Find route shares \(\vect{z}_k\) that approximate desired movement flows:
  \begin{equation}
    \min_{\vect{z}_k}
      \|W(A_k\vect{z}_k-\vect{f}_k^{\mathrm{target}})\|_2^2
      + \lambda_z\|\vect{z}_k-\vect{z}_k^{\mathrm{prev}}\|_2^2,
    \quad
    \vect{z}_k\geq0,
    \quad
    \vect{1}^{\trans}\vect{z}_k=1.
    \label{eq:path-decomposition}
  \end{equation}
  \item Assign recommendations by randomized or balanced rounding so aggregate route
  counts follow \(\vect{z}_k\).
\end{enumerate}

The candidate set limits action complexity but can exclude a socially useful route.
Report how often expanding \(K_{\mathrm{path}}\) changes the objective.

\subsection{Next-hop sampling method}

Alternatively, use the movement probabilities in \cref{eq:turn-softmax} to sample a
next hop as vehicles reach decision nodes.  This responds naturally to changing state
but needs loop prevention and can make a driver's displayed route unstable.  For a
navigation interface, commit each driver to a route for a minimum interval and reroute
only when expected improvement exceeds a threshold.

\subsection{Partial control and compliance}

The initial model can assume a controlled fraction \(\rho\) of drivers and route the
remainder according to UE-like or dynamic shortest-path behavior.  The simulator then
combines controlled and uncontrolled movement flows.  Different networks may require
very different controlled fractions to approach SO; optimized-fleet research has
demonstrated the importance of such penetration thresholds \citep{battifarano2023}.
Later work should learn compliance as a function of recommended detour and historical
trust rather than treating \(\rho\) as fixed \citep{jiang2024,yun2024}.

\section{Experimental protocol}

\subsection{Baselines}

Every experiment should include the following where applicable:

\begin{enumerate}[leftmargin=1.5em]
  \item free-flow shortest paths;
  \item dynamic shortest paths using current estimated travel time;
  \item static or dynamic user equilibrium;
  \item static marginal-cost assignment;
  \item receding-horizon DSO oracle;
  \item model-predictive control with the same CTM and shorter computational budget;
  \item classical pressure/backpressure routing or control where definitions align;
  \item a non-graph multilayer perceptron with comparable parameter count;
  \item a homogeneous GCN/GRU surrogate;
  \item the proposed heterograph model without projection;
  \item the proposed model without decision-focused training;
  \item an existing dynamic traffic-assignment graph model when a reproducible
  implementation is available \citep{liu2024hstgsn}.
\end{enumerate}

Max-pressure and backpressure methods are relevant stability-oriented comparators, but
their objective is not identical to minimum total travel time.  Classical max-pressure
analysis also commonly uses point queues with unlimited storage, which should be
acknowledged when comparing it with spillback-aware CTM control \citep{varaiya2013}.

\subsection{Experiment matrix}

\begin{table}[ht]
\centering
\caption{Minimum recommended experiment matrix.}
\label{tab:experiments}
\begin{tabularx}{\textwidth}{@{}p{0.09\textwidth}p{0.28\textwidth}X@{}}
\toprule
ID & Experiment & Question answered \\
\midrule
E0 & Braess and tiny networks & Are UE, SO, conservation, and price-of-anarchy calculations correct? \\
E1 & Seen topology, new demand & Can the model interpolate across demand profiles? \\
E2 & Seen topology, extreme demand & Does it remain useful near and beyond capacity? \\
E3 & Entirely unseen topology & Does it generalize structurally rather than memorize a graph? \\
E4 & Unseen incidents & How does performance degrade under capacity and reachability shifts? \\
E5 & CTM-to-SUMO transfer & Does improvement survive microscopic-model mismatch? \\
E6 & Controlled fraction sweep & How much participation is required for system benefit? \\
E7 & Network and commodity scaling & What are end-to-end runtime and memory growth? \\
\bottomrule
\end{tabularx}
\end{table}

SUMO is an appropriate microscopic evaluation environment and provides routing and
dynamic-assignment facilities \citep{lopez2018}.  It should be used as an independent
test environment after CTM training, not tuned until it agrees with the learned model.
Calibrate broad physical quantities such as free-flow time, capacity, and storage, then
report the remaining simulator gap.

\subsection{Primary metrics}

\paragraph{System performance.}
Report total system travel time, mean completed-trip time, median and 95th percentile,
throughput, unfinished trips, source-queue time, and maximum network accumulation.

\paragraph{Optimality.}
Report relative oracle gap

\begin{equation}
  g_J = \frac{J(\pi_\theta)-J(\pi^\star)}
              {J(\pi^\star)+\varepsilon},
\end{equation}

the recovered UE--SO gap \(1-\gap(\pi_\theta)\), and solver-bound-aware regret when the
oracle is not certified optimal.

\paragraph{Feasibility.}
Report conservation residual, capacity-violation magnitude, invalid-turn rate,
recirculating flow, demand completion, raw-to-projected action distance, and the number
of projection failures.

\paragraph{Robustness.}
Report performance as a function of demand shift, capacity loss, forecast error,
sensor noise, control delay, and controlled-vehicle fraction.

\paragraph{Computation.}
Report median, 95th, and 99th percentile wall-clock latency; hardware; batch size;
memory; graph preprocessing; projection; path decomposition; and end-to-end control
time.  Comparing only neural forward-pass time with total solver time would be
misleading.

\paragraph{Fairness and user burden.}
Even when fairness is not optimized, report the distribution of individual detour
ratios relative to selfish routing, worst affected OD groups, and the Gini coefficient
or another predeclared dispersion measure.  System-optimal guidance can improve the
aggregate while disadvantaging particular users \citep{jahn2005}.

\subsection{Statistical design}

Use common random numbers: every controller should face the same demand, incident, and
microscopic-driver seeds in a paired comparison.  Predeclare primary metrics and use
at least enough independent scenarios for interval estimates to stabilize; 30 paired
seeds per central condition is a reasonable initial target, not a universal rule.
Report paired mean or median differences with 95\% confidence intervals and effect
sizes.  Use bootstrap intervals if distributions are strongly skewed.  Treat time
steps within one simulation as dependent observations, not as thousands of independent
replicates.  If many secondary hypotheses are tested, distinguish exploratory results
or control the family-wise error rate.

\section{Implementation roadmap}

\subsection{Phase 0: correctness before learning}

\begin{enumerate}[leftmargin=1.5em]
  \item Implement static UE and SO on the Braess and Sioux Falls networks.
  \item Reproduce analytically checkable parallel-link and Braess examples.
  \item Verify \cref{eq:gap}: UE maps near one and SO maps near zero.
  \item Implement CTM conservation and test vehicle counts under no exits, known exits,
  and oversaturated source demand.
  \item Solve small DSO problems and compare the optimizer with exhaustive enumeration
  where possible.
\end{enumerate}

Do not begin GNN training until these tests pass.  A learned model can conceal errors in
the simulator or oracle by imitating them accurately.

\subsection{Phase 1: same-topology surrogate}

Train a destination-conditioned GCN/GRU on one network.  Use movement split logits,
masked softmax, and the exact CTM node allocator.  Demonstrate that it can imitate the
first oracle action and improve total travel time over dynamic shortest paths on held-out
demand profiles.  This phase validates the pipeline but is not yet the main novelty.

\subsection{Phase 2: heterogeneous topology-general model}

Add OD tokens, relation-specific message passing, physically normalized features, and
full topology splits.  Train on synthetic graph families and selected benchmark
networks.  Tune only on the validation topologies.  Freeze the experiment plan before
running the final test networks.

\subsection{Phase 3: projection and decision-focused training}

Add the projection in \cref{eq:projection}.  First train it as a postprocessor, then
differentiate through it if runtime and numerical conditioning permit.  Fine-tune with
short differentiable CTM rollouts and gradually lengthen the horizon.  Run ablations to
separate gains from architecture, projection, and decision loss.

\subsection{Phase 4: robustness and independent simulation}

Evaluate incidents, demand shift, noisy forecasts, and partial control.  Transfer the
frozen controller into SUMO through its traffic-control interface.  SUMO is microscopic
and continuous, whereas the training CTM is aggregate and discrete; performance loss is
expected and is itself an important result \citep{lopez2018}.

\subsection{Suggested software organization}

\begin{verbatim}
project/
  configs/             # frozen experiment and network configurations
  data_raw/            # immutable TNTP/OSM inputs and metadata
  data_processed/      # cell graphs and scenario tensors
  oracle/              # CTM, DSO formulation, solver adapters
  models/              # encoders, heterograph GNN, temporal module, decoders
  projection/          # feasibility layer and numerical tests
  routing/             # path generation and aggregate-to-driver assignment
  simulators/          # CTM rollout and SUMO interface
  experiments/         # training, evaluation, ablation entry points
  tests/               # conservation, reachability, oracle, and regression tests
  artifacts/           # versioned tables and figures generated from raw results
\end{verbatim}

Each saved result should include the source revision, configuration, network hash,
random seeds, solver version, solver tolerances, hardware, and model checkpoint hash.

\section{Failure modes and diagnostic tests}

\begin{longtable}{@{}p{0.22\textwidth}p{0.35\textwidth}p{0.36\textwidth}@{}}
\caption{Major failure modes and required diagnostics.}\label{tab:failures}\\
\toprule
Failure mode & Why it occurs & Diagnostic or mitigation \\
\midrule
\endfirsthead
\toprule
Failure mode & Why it occurs & Diagnostic or mitigation \\
\midrule
\endhead
Low flow error, poor travel time & Congestion cost is nonlinear near capacity. & Train and select on rollout objective; stratify error by volume/capacity ratio. \\
Feasible one-step action, unstable rollout & Small errors compound and shift visited states. & On-policy aggregation, longer rollouts, terminal penalties, stress tests. \\
Projection does all the work & Raw network outputs ignore physics. & Report raw violations and projection distance; penalize corrections. \\
Good random split, bad new graph & Scenarios from the same topology leaked into test. & Hold out complete networks before scenario generation. \\
Source holding artifact & Controller delays entry to reduce in-network occupancy. & Count source queues in TSTT and report admitted demand. \\
Route cycles & Local next-hop probabilities lack a global destination potential. & Reachability masks, loop penalty, destination potentials, loop-free path decomposition. \\
Oversmoothing & Deep message passing makes road embeddings indistinguishable. & Residual layers, limited depth, multiscale pooling, OD cross-attention. \\
Poor long-range OD reasoning & Origin and destination are many physical hops apart. & Virtual OD tokens/links and destination-conditioned decoding. \\
Oracle imitation ceiling & Oracle uses an inaccurate model or weak time limit. & Certified gaps, multiple solver budgets, SUMO transfer, lower-bound reporting. \\
Deployment latency hidden & Route generation or projection dominates forward pass. & Report complete pipeline latency and tail percentiles. \\
Unfair recommendations & SO sacrifices some users for aggregate benefit. & Report detour distributions; later add explicit constraints or rotation. \\
Control has no effect & Too few users accept recommendations. & Sweep controlled fraction and model compliance in extension experiments. \\
\bottomrule
\end{longtable}

\section{Ablation and claim discipline}

The final paper should align every claim with an experiment:

\begin{table}[ht]
\centering
\caption{Claims and minimum supporting evidence.}
\label{tab:claims}
\begin{tabularx}{\textwidth}{@{}p{0.29\textwidth}X@{}}
\toprule
Claim & Minimum evidence \\
\midrule
Near-system-optimal & Oracle-relative objective gaps with solver certification or bounds on all reported test scenarios. \\
Faster & End-to-end latency including preprocessing, projection, and route generation on identical hardware assumptions. \\
Physics-constrained & Hard feasibility or explicit residual distributions before and after projection. \\
Topology-generalizing & Entire road networks excluded from training and hyperparameter selection. \\
Dynamic & Closed-loop, time-dependent demand and queue propagation; not merely independent static snapshots. \\
Robust & Predeclared demand, incident, and observation shifts with degradation curves. \\
Scalable & Runtime and memory curves over network size and active commodities. \\
Improves price of anarchy & Same demand and feasible set for UE, learned policy, and SO; report \cref{eq:gap}. \\
\bottomrule
\end{tabularx}
\end{table}

Avoid saying the model ``finds the optimal route for all drivers.''  The system optimum
minimizes an aggregate objective and may lengthen some individual trips.  Say instead
that the model \emph{approximates a dynamic system-optimal aggregate routing policy}
under a stated demand, compliance, and traffic model.

\section{Proposed hypotheses and decision criteria}

The following hypotheses are specific enough to preregister:

\begin{description}[leftmargin=1.6in,style=nextline]
  \item[H1: decision quality]
  On held-out demand profiles for seen topologies, the projected GNN has lower TSTT than
  dynamic shortest-path and UE baselines and a small oracle-relative gap.
  \item[H2: topology transfer]
  On completely unseen topologies, physically normalized heterogeneous features and OD
  tokens reduce the oracle-relative gap compared with a homogeneous GCN baseline.
  \item[H3: feasibility]
  Projection eliminates operational constraint violations, while decision-focused
  training reduces the average projection correction relative to imitation alone.
  \item[H4: robustness]
  Incident-augmented training reduces performance degradation on unseen incidents
  without materially degrading nominal performance.
  \item[H5: computational value]
  The surrogate's end-to-end decision latency grows more slowly in practice than
  repeatedly solving the selected DSO oracle at matched solution quality.
\end{description}

Illustrative success criteria for continuing to a larger study are: (i) less than
5--10\% median oracle-relative TSTT gap on held-out scenarios, (ii) a substantial
reduction of the UE--SO gap, (iii) zero post-projection hard violations, (iv) useful
performance on at least one unseen topology, and (v) a clear end-to-end latency
advantage.  These thresholds should be revised based on pilot solver accuracy and
network difficulty before final preregistration.

\section{Extensions beyond the first paper}

\subsection{Lyapunov-shielded routing}

A follow-up can combine the surrogate with a stability-oriented controller.  For queue
state \(\vect{q}(t)\), define a Lyapunov function such as

\begin{equation}
  L(\vect{q}(t)) = \frac{1}{2}\sum_i w_i q_i(t)^2
\end{equation}

or a barrier-like spillback term that grows as occupancy approaches storage.  A safety
layer can restrict neural actions to those satisfying an upper bound on estimated
Lyapunov drift.  Backpressure and max-pressure traffic control provide relevant
stability foundations \citep{varaiya2013,gregoire2014,zaidi2016}.  The important new
question would be whether a learned correction improves delay without sacrificing a
stated stability region.

\subsection{Fairness and learned compliance}

Another extension treats the platform as a leader and drivers as probabilistic
followers.  Compliance may depend on recommended detour, trust, and incentives.  The
controller can minimize TSTT subject to per-user or per-group detour constraints.
Constrained system-optimal routing is an established fairness foundation
\citep{jahn2005}; the new work would need dynamic, learned behavior and long-run trust.

\subsection{Learned tolls or information design}

Rather than issuing prescriptive routes, the GNN can output link tolls, credits, or
information signals intended to induce socially beneficial equilibrium.  This is more
incentive-compatible but introduces a bilevel equilibrium problem and distributional
questions about who pays.  It should not be mixed into the initial surrogate paper.

\section{Conclusion}

A graph neural surrogate for DSO-DTA is technically plausible and potentially useful
because expensive optimization can be shifted offline while fast graph-based inference
is used online.  The project is publishable only if it goes beyond flow prediction.
The model must be evaluated as a constrained decision system: it should approximate
dynamic system-optimal actions, preserve vehicle conservation and capacity limits,
generalize to entirely unseen networks and disruptions, and demonstrate a true
end-to-end computational advantage.  A CTM oracle, heterogeneous OD-aware graph,
feasibility projection, decision-focused fine-tuning, and topology-held-out evaluation
form a coherent path to that goal.

\bibliographystyle{plainnat}
\bibliography{references}

\end{document}
